% Source: sources/returns/S6_04_MAPS_MANIFOLDS/sections/manifolds.tex:4904–5171
\paragraph{Theorem D.2.25 (Proved: all initial rows and all regular logarithms in the semistable sector).}
Consider any of the $62$ natural initial rows of D.2.22 on the four
all-multiplicative profiles $45,47,55,56$. Write $P=aG,Q=bG$, with
$a,b>0$. Use the primitive Mumford model at each multiplicative fibre.
An original zero-free filling with zero line multidegree and one
component orbit may be used in its place; it has the same marked
primitive rank-one model. Perform any finite number of regular
logarithmic modifications, with data
\[
 m_i>0,\qquad v_i=(x_i,y_i,z_i,n_i)\in\mathbb Z^4,
 \qquad\gcd(m_i,x_i,y_i,z_i,n_i)=1.                           \tag{D.SC1}
\]
The coordinates are $(e_1,e_2,\delta,\gamma)$. Put
$\omega=a e_1^*e_2^*+b\gamma^*\delta^*$. The initial rows have
$\gcd(a,b)=1$. Retain the peripheral integer $e$ of D.2.23;
one may keep every order-one site and use the canonical value $e=0$.

Here is a complete integral homology algorithm without a bound on
orders, number of sites, or translation coordinates. Let $S$ be the
Seifert relation matrix (D.FG2), and write
\[
 \operatorname{coker}S=\mathbb Z^r\oplus A_1,\qquad r\in\{0,1\}.
\]
Form an integral presentation $\mathcal C$ with generators
$s$ and $t_i(u)$, linear in $u\in\Lambda=\mathbb Z^4$, and relations
\[
 \begin{split}
 m_i t_i(u)&=\omega(u,v_i)s\quad(u=e_1,e_2,\delta,\gamma),\\
 t_i(v_i)&=0,\\
 \sum_i t_i(\delta)+be\,s&=0.
 \end{split}                                                \tag{D.SC2}
\]
Write $T=\operatorname{Tor}\operatorname{coker}\mathcal C$.
Then $\operatorname{coker}\mathcal C$ has free rank $r$ and
\[
\begin{array}{c|c}
 j&H_j(X;\mathbb Z)\\\hline
 0,6&\mathbb Z\\
 1&\mathbb Z^r\oplus A_1\\
 2,3&\mathbb Z^{2r}\oplus T\\
 4&\mathbb Z^{2r}\oplus A_1\\
 5&\mathbb Z^r.
\end{array}                                                 \tag{D.SC3}
\]
The fundamental group is exactly (D.FG1), rather than merely the
abelian group in degree one. Smith reduction of the explicit matrix
(D.SC2) computes $T$, including all extensions.

If this group is trivial, let
\[
 M=\prod_i m_i,\quad
 X_0=\sum_i(M/m_i)x_i,\quad Y_0=\sum_i(M/m_i)y_i,\quad
 g=\gcd(b,M,X_0,Y_0).
\]
For no displayed sites use $M=1$ and $X_0=Y_0=0$.
Then
\[
 H_2(X)=H_3(X)=
 \mathbb Z/g\oplus\mathbb Z/M\oplus\mathbb Z/(bM/g).          \tag{D.SC4}
\]
Consequently a sphere occurs precisely when the group criterion
(D.FG3) holds and $b=M=1$. In the canonical marking with all sites
retained, this says that every regular logarithmic order is one and
$|\sum_i n_i|=1$.

There are exactly $17$ permitted initial rows with $b=1$:
\[
\begin{array}{c|c|c|c}
 \text{profile}&\text{zero fibre}&(a,b)&\text{marked positions}\\\hline
 45&I_1&(8,1)&4\\
 47&I_2&(7,1)&1\\
 47&I_1&(14,1)&3\\
 55&I_4&(5,1)&1\\
 55&I_1&(20,1)&3\\
 56&I_5&(6,1)&1\\
 56&I_3&(10,1)&1\\
 56&I_2&(15,1)&1\\
 56&I_1&(30,1)&2.
\end{array}                                                 \tag{D.SC5}
\]
This is the complete natural sphere table for these four profiles,
including original-versus-Mumford choices: any additional original
cap not covered by the first paragraph is excluded by the positive
local defect in D.2.22.

More generally, every simply connected output in these $17$ rows,
with arbitrary data (D.SC1), is oriented diffeomorphic to
\[
                              \widetilde\Sigma_0N_M,         \tag{D.SC6}
\]
with the twisted framing and the Smale five-manifold of D.2.14.
Thus its full smooth type depends only on $M$, including when $M$
is even. For $b>1$, (D.SC2)--(D.SC4) give the complete integral
homology and exclude a sphere; no smooth classification of all
those torsion types is asserted here.

\begin{proof}
\emph{1. The actual marked caps and the constant integral calculation.}
At the component zero the valuation lattice has index one by D.2.22,
so its cap is the primitive hexagonal model. Every other
multiplicative monodromy in these $62$ rows has rank-one difference
and Smith width one, as checked in D.2.23. Its local Mumford cap is
the marked one-node model times $T^2$. The integral specialization
onto its $H_2$, and that of the hexagonal cap, are onto; their
boundary kernels are exactly the normal tori swept by
$\operatorname{Fix}(A_i)$. The toric angular coordinates give these
as pure swept classes, with zero fibre-area term. The same statement
for an original zero-degree, one-orbit cap follows by quotienting
its component cycle by $Q$, using the unit-factor isotopy of D.2.12.

These normal markings agree with the section-line construction.
Rigidify $M$ along $O$ on the punctured base. Put the divisor of this
rigidification at the chosen zero; its integral change there is
multiplication of the line coordinate by $t^h$. On the infinite
toric model this translates the fan by an integral lattice vector,
and hence extends over the cap. At all other points the
rigidification is a unit and the earlier relative unit isotopy
applies. Thus this normalization does not assume that equality of
vertical forms alone determines mixed swept classes.

The exact program \texttt{manifolds\_semistable\_catalogue.py}
performs the graph and exterior calculations for all $62$ rows.
The degree-two fibre coinvariants are freely generated by $s$,
with primitive detector $\omega$; all its Smith factors are one.
Adjoin one regular puncture, with monodromy the identity. In the
notation of D.2.13, the $14$ old normal columns form a primitive
rank-$14$ sublattice of the rank-$17$ graph kernel. Their quotient
is freely generated by the regular sweeps of $e_1,e_2,\gamma$,
while the sweep of $\delta$ is zero. All $62$ left inverses and
quotient matrices are saved, so this is an exact integral
certificate, not a numerical rank inference.

Adding another regular puncture adds four free graph coordinates.
The old normal lattice remains primitive by projection onto the
previous graph. Since every monodromy fixes $\delta$, the peripheral
Fox maps satisfy $\sum j_i(\delta)=0$. Induction therefore gives
all regular sweeps freely, with the sole relation
$\sum_i t_i(\delta)=0$. The degree-two exterior summand remains
$\mathbb Zs$. A change of global peripheral lift by $e\gamma$
adds $be\,s$ to this relation: the correction to a swept $\delta$
cycle is its wedge with that translation, and
$\omega(\delta,\gamma)=-b$. Equivalently, under
$v_i\mapsto v_i+m_i h_i\gamma$ the new sweep is
$t_i'(u)=t_i(u)+\omega(u,h_i\gamma)s$;
$e\mapsto e+\sum h_i$ makes (D.SC2) invariant. This checks its
sign against (D.FG1).

\emph{2. The complete regular-log kernel.}
At a regular log point its reduced torus has lattice
$\Lambda'=\Lambda+\mathbb Z(v_i/m_i)$. The nearby fibre inclusion
and the base circle send a boundary bivector respectively to
$z$ and $u\wedge(v_i/m_i)$ in $\bigwedge^2\Lambda'$.
This map is onto. The condition in (D.SC1) means $v_i$ is primitive
modulo $m_i$; the integral normal lattice is exactly
$m_i\Lambda+\mathbb Z v_i$. This can be checked in a basis in which
$v_i$ modulo $m_i$ is the first basis vector, using the Euclidean
algorithm. Its generators give the full normal columns
\[
 (-\omega(u,v_i),\ m_it_i(u)),\qquad (0,t_i(v_i)).
\]
They are precisely the first two lines of (D.SC2). The argument
identifies the full integral kernel, including its index, not only
its rational span. Since specialization is onto every cap in degree
two, Mayer--Vietoris reduces to this cokernel.

The remaining first MV kernel is free: all boundary $H_1$ groups
are free in these primitive toric models. Its rank is $r$, because
the domain and target of the first boundary map have equal rational
dimension and its cokernel is $H_1(X)$. Hence it contributes a
split $\mathbb Z^r$ to $H_2(X)$. Over $\mathbb Q$, eliminate the
sweeps in (D.SC2); its sole residual relation is
\[
                 b\left(e-\sum_i n_i/m_i\right)s=0.
\]
Thus its free rank is the same $r$ as the Seifert matrix.
Only the primitive hexagonal cap contributes Euler characteristic,
so $\chi(X)=2$. Integral Poincare duality and the universal coefficient
theorem give (D.SC3), including the identical torsion in $H_2,H_3$.
This also proves the formula when the rational Euler number is zero;
that case is not obtained by deleting a nonzero-order torsion factor.

For one site in canonical marking, quotienting the zero $\delta$
sweep displays the concrete four-by-five matrix
\[
 \begin{pmatrix}
 -ay&ax&bn&-bz&0\\
 m&0&0&0&x\\0&m&0&0&y\\0&0&0&m&n
 \end{pmatrix}.                                             \tag{D.SC7}
\]
Its maximal minors are, up to signs,
$bm^2n\,(m,n,z,y,x)$. In particular, whenever $n\ne0$,
the finite degree-two cokernel has order $bm^2|n|$.
The implementation accepts the general multi-site matrix (D.SC2),
and verifies its comparison with all constant graph presentations.

\emph{3. The simply connected Smith form.}
Assume (D.FG1) is trivial. After removing order-one sites there are
at most two multiplicities greater than one, they are coprime, and
every $n_i$ is a unit modulo its $m_i$. Put
\[
                  D=M\left(\sum_i n_i/m_i-e\right).
\]
The group criterion gives $D=\pm1$. Work over $\mathbb Z_p$.
If $p$ does not divide $M$, all sweeps can be eliminated and the
cokernel is $\mathbb Z_p/(b)$. If $p\mid M$, it divides exactly
one multiplicity, say $m_i$. For every other site $m_j$ is a unit,
so its $\delta$ sweep equals $-bn_j/m_j$ times $s$. The global
relation then eliminates $t_i(\delta)$. Its own normal relation
becomes $bD/(M/m_i)$ times $s$, and hence is equivalent to $bs=0$.
Since $n_i$ is a unit, its relation $t_i(v_i)=0$ eliminates
$t_i(\gamma)$, whose remaining normal relation is redundant by
alternation of $\omega$. The three remaining generators have
presentation
\[
 \begin{pmatrix}
 b&0&0\\ -ay_i&m_i&0\\ ax_i&0&m_i
 \end{pmatrix}.                                             \tag{D.SC8}
\]
Its determinantal divisors are $g_i,m_i g_i,bm_i^2$, where
$g_i=\gcd(b,m_i,x_i,y_i)$ in $\mathbb Z_p$; we used
$\gcd(a,b)=1$. The Chinese remainder theorem packages all these
primary invariant factors as (D.SC4), because
$X_0,Y_0$ reduce at $p\mid m_i$ to unit multiples of $x_i,y_i$.
The invariant factors divide in the stated order, since $g$ divides
both $b$ and $M$. This proves the formula for every integer twist,
not by extrapolating from the finite checks.

The group in (D.SC4) vanishes exactly for $b=M=1$. Simple connectivity,
Euler number two, and integral duality then give a homotopy sphere;
$\Theta_6=0$ identifies it with the standard $S^6$. The seventeen
initial rows in (D.SC5) are the complete $b=1$ filter of D.2.22.
An excluded original cap has nonzero Chern degree or multiple
component orbits and would already force $b_2>0$, so cannot supply
additional sphere rows.

\emph{4. The smooth type when $b=1$.}
We verify the hypotheses of the geometric proof D.2.14 rather than
inferring a smooth type from homology. At every zero-free
multiplicative point the component equation forces $P$ to be narrow
when $Q=G$; the component action of $Q$ is primitive and transitive. The code
checks these component statements and, more directly, that adjoining
$\delta$ to the actual primitive vanishing vector gives two unit
Smith factors at every such site. Thus $\delta$ stays primitive
modulo the vanishing line and its circle is free on the cap; width
one alone is not being used to infer this. At the zero, index one makes the valuation lattice the full original
short lattice, in which $\delta$ is the second coordinate. The code
checks its primitive inclusion in that lattice in all seventeen rows.
The actual hexagonal fan weights are therefore the same $(-1,0,1)$ as in D.2.14, with one fixed rational curve and
normal bundle $\mathcal O(-1)_{+1}\oplus\mathcal O(-1)_{-1}$.
For every regular log cap, $\gcd(m_i,n_i)=1$ implies that a scalar
angle lying in $\Lambda+\mathbb Z(v_i/m_i)$ is integral. Hence no
new circle stabilizers occur.

The fixed curve lies in the cyclic fibre summand generated by $s$.
The $\delta$ normal relations together with the global relation,
when $b=1$, are the unimodular Seifert matrix because $D=\pm1$.
They kill $s$ integrally. The fixed sphere is therefore nullhomologous
and nullhomotopic. Its equivariant surgery gives
$Y\cong X\#(S^3\times S^3)$ with a free, projectably spin circle
action, exactly as in D.2.14. The same orientation convention in its
normal-coordinate conjugation is retained. Its spin simply connected
quotient has $H_2=\mathbb Z\oplus(\mathbb Z/M)^2$, with primitive
Euler class. Smale--Barden and the spin branch of
\cite[Theorem A]{manGR25} give
$Y\cong(S^3\times S^3)\#\widetilde\Sigma_0N_M$.
Cancellation by the full Zhubr invariants, as justified in D.2.14,
proves (D.SC6), including the two-primary refined invariant.
\end{proof}
